\subsection{Refining the Scale by Bisection}

\begin{derivationbox}[Constructing finer marks]
The compass-and-straightedge construction does not stop at integer marks.
Bisecting each unit segment constructs halves. Bisecting again constructs
quarters, then eighths, and so on. Repeating this process on every integer
interval and on both sides of the origin gives all dyadic marks
\[
\boxed{\frac{k}{2^n}},
\qquad k\in\mathbb Z,
\quad n\in\mathbb N.
\]
\end{derivationbox}

\begin{center}
\begin{tikzpicture}[scale=1.1]
    \draw[axis] (-0.4,0) -- (4.4,0);
    \foreach \x/\lab in {0/0,1/{\frac14},2/{\frac12},3/{\frac34},4/1} {
        \draw (\x,0.1) -- (\x,-0.1) node[below=4pt] {$\lab$};
    }
    \node[above] at (2,0.35) {successive bisections of one unit};
\end{tikzpicture}
\end{center}

\begin{theorembox}
The dyadic marks are dense on the axis: every nonempty interval contains a
number of the form $k/2^n$. They provide arbitrarily fine numerical
approximations to positions on the line.
\end{theorembox}

\begin{definitionbox}
\textbf{Continuity assumption.} Once an origin, a positive direction, and a unit
length are fixed on a Euclidean line, every point of the line corresponds to
exactly one real number, and every real number corresponds to exactly one point
of the line.
\end{definitionbox}

\begin{interpretationbox}[What continuity contributes]
Finite straightedge-and-compass operations construct many particular positions,
including every dyadic mark, but not every real position. Continuity supplies
the complete real scale. After this additional assumption, every point of each
axis has a unique signed real label.
\end{interpretationbox}
